% !TEX program = xelatex
\documentclass[UTF8,11pt]{ctexart}

\usepackage[a4paper,margin=2.2cm]{geometry}
\usepackage{hyperref}
\usepackage{bookmark}
\usepackage{listings}
\usepackage{xcolor}
\usepackage{enumitem}
\usepackage{graphicx}
\usepackage{tabularx}
\usepackage{amsmath}

\hypersetup{
  colorlinks=true,
  linkcolor=blue,
  urlcolor=blue,
  citecolor=blue
}

\lstset{
  basicstyle=\ttfamily\small,
  breaklines=true,
  columns=fullflexible,
  frame=single,
  rulecolor=\color{black!20},
  backgroundcolor=\color{black!2}
}

\title{Codex \,+\,Claude Code MCP：上下文注入与跨会话记忆（第一性原理方案）}
\author{jqwang \quad (draft)}
\date{\today}

\begin{document}
\maketitle

\begin{abstract}
当前我们在 Codex 中集成了一个 \texttt{claude\_code} 的 MCP 工具：当提示词中“呼唤 Claude 做事情”时，Codex 会转发调用到 Claude CLI 侧执行。但 MCP 调用默认只携带该次工具参数（prompt/workFolder 等），\textbf{不会自动携带当前 session 的历史上下文}，因此 Claude 很容易在缺少前情的情况下做出偏差实现。

本文从第一性原理出发，给出一套“\textbf{调用时上下文注入（Context Packet）} + \textbf{跨会话持久记忆（SQLite thread\_memory）}”的组合方案，并说明其与 Codex 官方方向（\texttt{get\_memory} 工具 + SQLite）如何对齐，以及如何验证/迭代。本文同时记录截至 \textbf{2026-02-09} 的落地实现现状与后续里程碑。
\end{abstract}

\tableofcontents
\newpage

\section{问题陈述}
\subsection{现象}
\begin{itemize}[nosep]
  \item 我们已能通过 MCP 调用 Claude Code 执行任务；
  \item 但当用户在同一 session 中多轮讨论后再次呼唤 Claude，Claude 侧常常\textbf{不知道之前的决策、约束、上下文摘要}；
  \item 结果：实现偏差（例如不遵守之前约定、忽略已有结论、重复工作）。
\end{itemize}

\subsection{根因（结构性）}
MCP 调用本质上是一次独立 RPC：Codex 会把工具参数 JSON 原样转发到 MCP server。
\textbf{除非我们显式把“历史摘要/上下文”作为参数传过去，否则 Claude CLI 侧拿不到任何 session 历史。}

\section{目标与非目标}
\subsection{目标（Goals）}
\begin{enumerate}[nosep]
  \item Claude Code 在被调用时，能获得\textbf{足够的高信噪比上下文}：
    \begin{itemize}[nosep]
      \item 当前工作目录（影响 \texttt{CLAUDE.md} 等本地指令文件的生效）；
      \item 当前 session 的摘要（不含推理链，只含结论/约束/决定）；
      \item 最近若干轮对话摘录（用于任务细节与引用）。
    \end{itemize}
  \item 在 Codex 内部建立\textbf{跨会话}可复用的“线程记忆”机制：同一 thread 被 resume 后仍可取回摘要。
  \item 方案与 Codex 官方方向对齐：SQLite \texttt{thread\_memory} + \texttt{get\_memory} 工具。
\end{enumerate}

\subsection{非目标（Non-goals）}
\begin{itemize}[nosep]
  \item 不把完整 transcript 全量塞给 Claude（成本/泄露风险/长度不可控）。
  \item 不传“思考过程/推理链”（只传总结与必要上下文）。
  \item 不要求 Claude Code 直接读取 Codex 的私有 state DB（保持边界清晰）。
\end{itemize}

\section{第一性原理：为什么需要\,``Context Packet}
\subsection{语言模型执行任务的输入假设}
将模型视为函数：
\[
  y = f(\text{instructions}, \text{context}, \text{task})
\]
若 \texttt{context} 缺失，模型会用默认先验填补空白（hallucination / wrong assumptions）。

\subsection{高信噪比上下文的最小闭包}
对“代写/代做代码”类任务，最小闭包通常包含：
\begin{itemize}[nosep]
  \item \textbf{约束与偏好}：例如编码规范、禁止改动的区域、测试要求；
  \item \textbf{已做的结论/决策}：例如“我们用 resume 实现历史进入”、“不注入 reasoning trace”；
  \item \textbf{当前工作区环境}：cwd、构建工具、可用命令；
  \item \textbf{最近的具体需求}：用户刚刚提出的目标与验收标准。
\end{itemize}

\subsection{为什么不能直接传全量 transcript}
\begin{itemize}[nosep]
  \item 长度不可控，容易超上下文窗；
  \item 噪声大（工具输出、日志、反复讨论）；
  \item 风险高（可能包含密钥/隐私）；
  \item 成本高（token + 运行时间）。
\end{itemize}
因此，需要一份可控的、结构化的 ``Context Packet。

\section{Codex 现状（可复用的“上下文来源”）}
\subsection{Session 内：History / Compaction}
\begin{itemize}[nosep]
  \item Codex 内部可获取完整 transcript：\texttt{Session::clone\_history()}\, + \texttt{ContextManager::raw\_items()}。
  \item Codex compaction 会在 transcript 中写入一个 ``summary message（用户 role 的一段摘要前缀）。
  \item 这提供了\textbf{天然的“session 内摘要来源”}，且不包含模型推理链。
\end{itemize}

\subsection{跨会话：SQLite thread\_memory + get\_memory}
\begin{itemize}[nosep]
  \item SQLite migration：\texttt{thread\_memory(thread\_id, trace\_summary, memory\_summary, updated\_at)}。
  \item Tool：\texttt{get\_memory(memory\_id = thread\_id)} 读取该结构化 JSON（feature-gated：\texttt{memory\_tool}）。
  \item \textbf{已补齐关键链路：}实现了 thread\_memory 的自动生成/更新流水线（turn 完成 + compaction 后写入）。
\end{itemize}

\section{方案总览}
\subsection{两层结构}
\begin{enumerate}
  \item \textbf{调用时上下文注入（Context Packet）}：
    \begin{itemize}[nosep]
      \item 发生在每次调用 \texttt{claude\_code} MCP 工具时；
      \item 组装并注入一个 \texttt{context} 字段（字符串），且补齐 \texttt{workFolder}。
    \end{itemize}
  \item \textbf{跨会话持久记忆（Thread Memory）}：
    \begin{itemize}[nosep]
      \item 持久化在 SQLite；
      \item 可在 resume 后取回；
      \item 作为 Context Packet 的一部分（如果存在则拼进去）。
    \end{itemize}
\end{enumerate}

\subsection{Context Packet 建议结构}
\begin{lstlisting}
Working directory: /path/to/repo

Saved thread memory:
Trace summary:
...

Memory summary:
...

Session summary:
...

Recent chat excerpt:
User: ...

Assistant: ...
\end{lstlisting}

\section{Claude Code MCP：注入点与关键行为}
\subsection{注入点}
建议在 Codex 的 MCP 调用入口处，对 \texttt{tool == "claude\_code"} 做参数增强：
\begin{itemize}[nosep]
  \item 若未提供 \texttt{workFolder}：设置为 \texttt{turn\_context.cwd}。
  \item 若未提供 \texttt{context} 或为空：构建 Context Packet 并填入。
\end{itemize}

\subsection{为什么 workFolder 很关键（CLAUDE.md 生效）}
Claude Code 通常会读取当前工作目录（或指定 work folder）中的 \texttt{CLAUDE.md} 作为局部指令。
因此：\textbf{必须确保 MCP 工具默认在 Codex 的 session cwd 中运行}，否则 Claude 侧会漏掉本地约束。

\subsection{上下文抽取策略（不含推理链）}
\begin{itemize}[nosep]
  \item 优先使用 compaction summary message（如果存在）；
  \item 仅摘取 user/assistant 的文字消息；
  \item 限制最近消息条数与字节预算（避免过长）；
  \item 如存在 SQLite thread\_memory，拼接其 trace\_summary / memory\_summary（同样做截断）。
\end{itemize}

\section{Thread Memory：如何对齐官方方向}
\subsection{数据库与工具已经具备}
\begin{itemize}[nosep]
  \item 表结构与读工具已存在：\texttt{get\_memory}。
  \item 需要解决的关键问题是：\textbf{何时写入}与\textbf{如何生成}；当前已实现（见下节）。
\end{itemize}

\subsection{生成策略（已实现）}
\begin{enumerate}[nosep]
  \item 触发时机：
    \begin{itemize}[nosep]
      \item Turn 完成后（带 10 分钟节流），后台异步生成；
      \item Compaction 完成后，后台异步生成；
      \item 未来可补充：手动 backfill 命令，从 rollout 重新构建（用于历史 thread 补齐）。
    \end{itemize}
  \item 输入来源：rollout（.jsonl）或 normalized trace items。
  \item 生成方式：调用官方 \texttt{/v1/memories/trace\_summarize}（复用 Codex client wiring）。
  \item 写入：\texttt{upsert\_thread\_memory(thread\_id, trace\_summary, memory\_summary)}。
\end{enumerate}

\subsection{实现要点（当前代码行为）}
\begin{itemize}[nosep]
  \item feature gate：仅在启用 SQLite state DB 且启用 \texttt{Feature::MemoryTool} 时运行（避免在无 DB/无工具时引入副作用）。
  \item 内容过滤：memory trace 会排除 \texttt{FunctionCallOutput/CustomToolCallOutput/Reasoning/GhostSnapshot} 等高噪声项，仅保留 user/assistant 文字与工具调用元信息（不含工具输出）。
  \item 回退策略：在 trace summarize 失败/不支持时，优先写入 compaction summary 或最后一条 agent message（只写 summary，不写全量）。
  \item 写入时序：为避免 SQLite 中 thread 元数据尚未写入导致 upsert 失败，写入前会短暂轮询 thread 元数据可用性。
  \item 预算控制：trace item 数量与字节数上限（避免超长与高成本）。
\end{itemize}

\section{实现进展（截至 2026-02-09）}
\begin{enumerate}[nosep]
  \item SessionBar / Resume：Enter 触发真正的 resume 并回放历史消息，修复“进入历史 session 但无内容”的问题。
  \item Claude MCP Context Packet：
    \begin{itemize}[nosep]
      \item \texttt{workFolder} 缺省注入为 session cwd（使本地 \texttt{CLAUDE.md} 有机会生效）；
      \item \texttt{context} 缺省自动注入（Saved thread memory + Session summary + 最近摘录 + 同 cwd 的 project memories）。
    \end{itemize}
  \item Thread Memory 自动写入链路：turn 完成 + compaction 后写入 thread\_memory，并通过 \texttt{get\_memory} 可读回。
  \item Context Packet Builder：抽成可复用组件，并复用于 \texttt{spawn\_agent} 的初始消息（让 Codex 内部 sub-agent 也能拿到同一份上下文/记忆摘要）。
\end{enumerate}

\subsection{内容约束（记忆的“正确粒度”）}
\begin{itemize}[nosep]
  \item 只保留长期有价值的信息：项目背景、长期约束、关键架构决策；
  \item 不存短期任务细节（短期细节留给“最近消息摘录”）；
  \item 不存密钥/隐私；
  \item 可考虑按 cwd 关联（已有 API：last N memories for cwd）。
\end{itemize}

\section{Sessionbar / Resume：进入历史但无内容的问题}
问题的第一性原理解释：
\begin{itemize}[nosep]
  \item ``打开历史会话 本质是 \textbf{resume}：需要从 rollout 重建 thread 的 history 并恢复 session 配置；
  \item 若只是在 UI 层切换选中态而未触发 resume，新的 thread 并未加载历史，因此显示为空。
\end{itemize}

解决思路：把 SessionBar 的 Enter 动作映射到 ``ResumeSession(path)，由 App 统一走 \texttt{resume\_thread\_from\_rollout}。

\section{安全与隐私}
\begin{itemize}[nosep]
  \item Context Packet 应避免包含：工具原始输出、环境变量值、密钥文件内容。
  \item 必须做长度预算与截断。
  \item 未来可选：对历史做更严格的“敏感信息过滤”（例如 key patterns）。
\end{itemize}

\section{测试与验收}
\subsection{单元测试}
\begin{itemize}[nosep]
  \item context 抽取：summary 优先、只取 summary 之后的消息、截断行为。
\end{itemize}

\subsection{集成测试}
\begin{itemize}[nosep]
  \item TUI：SessionBar Enter 选择历史 session 触发 ResumeSession。
  \item MCP：\texttt{claude\_code} 工具调用参数中包含 workFolder/context（可通过 mock MCP server 验证）。
\end{itemize}

\section{里程碑（建议）}
\begin{enumerate}[nosep]
  \item M0（已完成）：调用时注入 workFolder + context（session 摘要 + 最近消息摘录）。
  \item M1（已完成）：接入 SQLite thread\_memory（如果存在则注入）。
  \item M2（已完成）：自动 summarization 并 upsert thread\_memory（官方 trace\_summarize + 回退策略）。
  \item M3（部分完成）：按 cwd 注入 project memories（同项目最近 N 条 thread\_memory 的 memory\_summary）；后续可做更智能的融合/去重/挑选策略。
  \item M4（建议）：提供 backfill / rebuild 命令：从 rollout 重建并写入 thread\_memory（便于对历史 thread 补齐）。
  \item M5（部分完成）：抽象可复用的 Context Packet Builder，并复用于 \texttt{claude\_code} 与 \texttt{spawn\_agent}；后续扩展到 Grok/Gemini/Review 等子 agent（多模型协作的一致性核心）。
\end{enumerate}

\end{document}
